\section{Individual Fairness Formulation}
\label{sec:individual-fairness}

\draftnote{DRAFT. Carried over from an earlier draft's Problem
Formulation section, restated as a proper (differentiable, directly
minimizable) loss rather than an indicator-weighted constraint, per
editorial instruction. No numbers from that earlier draft's own
experiments are carried over here or anywhere else in this paper.}

Let $u_i$ and $u_j$ denote two users with original mobility
trajectories $\tau_i$ and $\tau_j$, and let $\mathrm{sim}(\tau_i,
\tau_j)$ be a similarity function over trajectories, where higher
values indicate more similar mobility behavior. Let $P(\tau)$ denote a
privacy-gain measure associated with a trajectory $\tau$ -- in this
paper, $P$ is instantiated with the decile-level reduction statistics
of Section~\ref{sec:decile-reduction-metrics}, not the per-user ratio
an earlier formulation of this idea used (Section~\ref{sec:methods-changed}).
Individual fairness, following the general principle that similar
individuals should receive similar
treatment~\cite{dwork2012fairness}, requires that users with similar
trajectories receive comparable privacy gain. An earlier formulation
of this requirement stated it as a constraint, enforced only above a
similarity threshold $\delta$ and tolerating a fixed disparity budget
$\epsilon$:
\begin{equation}
\sum_{i,j} \mathbb{I}\big(\mathrm{sim}(\tau_i,\tau_j) \ge \delta\big)\,
\big|P(\tau_i) - P(\tau_j)\big| \le \epsilon.
\label{eq:IF-constraint}
\end{equation}
Equation~\ref{eq:IF-constraint} is not itself a loss that a gradient-based
optimizer can minimize: it is a single global inequality, not a
per-sample or per-pair quantity with a gradient, and the indicator
$\mathbb{I}(\cdot)$ has zero gradient everywhere it is differentiable
at all. Restating it as a proper loss means turning the constraint
into a weighted sum that is directly minimizable and whose minimum
(zero) exactly recovers the constraint being satisfied with equality
at $\epsilon=0$:
\begin{equation}
\mathcal{L}_{\mathrm{IF}} = \sum_{(i,j)} w_{ij}\,\big|P(\tau_i) - P(\tau_j)\big|,
\qquad
w_{ij} = \mathbb{I}\big(\mathrm{sim}(\tau_i,\tau_j) \ge \delta\big),
\label{eq:IF-loss}
\end{equation}
which is well-defined, differentiable in $P$ wherever $P$ itself is,
and minimized (at zero) exactly when every similar pair has identical
privacy gain. The similarity threshold $\delta$ and tolerance
$\epsilon$ move from being load-bearing parts of the training
objective to being evaluation-time choices: $\delta$ selects which
pairs $w_{ij}$ weights into the loss, and $\epsilon$ becomes a
reporting threshold for \emph{how fair is fair enough}, not a term the
optimizer itself needs. A hard threshold weight $w_{ij}\in\{0,1\}$ can
equivalently be replaced by a smooth kernel in $\mathrm{sim}(\tau_i,\tau_j)$
-- the Lipschitz-soft formulation evaluated in
Section~\ref{sec:phase4} does exactly this.

\subsection{Entropy as a similarity measure}
To instantiate $\mathrm{sim}(\tau_i,\tau_j)$, this paper (following the
same choice in the formulation it builds on) uses Shannon entropy
$H(\mathcal{T}_u)$ of a user's visitation distribution over a
discretized spatial grid as a proxy for mobility regularity: lower
entropy indicates a more routine, more predictable trace, which is
empirically also the easier trace to re-identify
(Section~\ref{sec:disparate-risk}). Entropy is attractive as a
similarity basis for exactly this reason -- it is computed from
exactly the same quantity individual fairness is meant to equalize
protection over -- and is computationally lightweight (a single-pass
grid-based statistic) and robust to small spatial or temporal
perturbations. This paper groups users into entropy \emph{deciles}
(Section~\ref{sec:evaluation-protocol}) rather than computing
$\mathrm{sim}(\tau_i,\tau_j)$ pairwise for every pair, which is what the
decile-level reduction and residual-risk-disparity statistics of
Section~\ref{sec:decile-reduction-metrics} report in place of a raw
pairwise fairness loss value.
